\documentclass[12pt]{article}
\usepackage[margin=1in]{geometry}
\usepackage{setspace}
\usepackage{parskip}
\usepackage{hyperref}
\usepackage{apacite}   % for APA citations; compile with BibTeX

\onehalfspacing

\title{Deconstructing My Stereotypes About ``Problem Students''\\
       at Pine Brook Elementary}
\author{Piter Garc\'ia}
\date{\today}

\begin{document}

\maketitle

% -----------------------------------------------------------------------
% SECTION 1 — Introduction: Naming My Bias
% -----------------------------------------------------------------------
\section{Introduction: Naming My Bias}

When I first arrived at Pine Brook Elementary as a student teacher, I
carried with me a set of ready-made ideas about which students would be
``problems.''  I expected certain children---usually the ones whose names
surfaced repeatedly in team meetings or behavior charts---to be
disruptive, unmotivated, or ``not ready'' for more complex,
technology-rich learning.  As a Dominican, immigrant, neurodivergent
educator-in-training entering a predominantly white suburban school, I
was also navigating my own position as someone often stereotyped and
misread in academic spaces.  Those overlapping identities shaped both how
I saw my students and how I anticipated they might see me.

Over the course of my placement, I began to notice that some of the
students most consistently described as ``behavior issues'' responded
differently with me than with other adults in the building.  In
transcripts of my interactions, I could see them taking initiative,
supporting peers, and engaging deeply with tools like AI, Minecraft,
Sphero robots, and circuitry kits when the task felt meaningful and
relational.  This paper uses those moments to examine how my prejudices,
biases, and stereotypes about so-called ``problem students'' were formed,
how they were reinforced by the institutional culture around me, and how
they were disrupted as I actually worked alongside these children.
Drawing on my field transcripts and course readings, I address four
guiding questions: how I connected thoughtfully with this group; what
myths and assumptions operated around them; what assets and funds of
knowledge I observed; and what I am committed to doing differently in my
future practice.

% -----------------------------------------------------------------------
% SECTION 2 — Connecting Thoughtfully With Labeled Students
% -----------------------------------------------------------------------
\section{Connecting Thoughtfully With Labeled Students}

% TODO: Describe how these students were first introduced to you
%       (meetings, documentation, informal comments, "frequent flyers").
% TODO: Explain specific steps taken to connect:
%       interests, predictable routines, tech-rich tasks, positioning as helpers/experts.
% TODO: Reflect on navigating your own safety/well-being as an evaluated,
%       minoritized student teacher while building those relationships.
% TODO: Note the tension of being "the tech person" vs. building trust
%       without undermining colleagues.
% TODO: Insert 1-2 transcript excerpts showing early interactions and
%       shifts in student responses over time.

\section{Myths, Prejudices, and Stereotypes Versus Transcript Reality}

% TODO: Name YOUR stereotypes at the start
%       ("they don't care," "they're just defiant," "can't handle advanced tools").
% TODO: Name SCHOOL stereotypes as evidenced in transcripts
%       (deficit language, low expectations, pathologizing of families).
% TODO: Place those myths next to 2-4 specific transcript moments where
%       the same students contradicted the labels when engaged and respected.
% TODO: Make explicit how those moments challenged your assumptions.

% -----------------------------------------------------------------------
% SECTION 4 — Assets and Funds of Knowledge I Observed
% -----------------------------------------------------------------------
\section{Assets and Funds of Knowledge I Observed}

% TODO: Reframe students through their strengths:
%       creativity, leadership, tech fluency, humor, care, persistence.
% TODO: Ground each asset in a specific transcript scene
%       (e.g., student troubleshooting devices; "disruptive" student
%        mastering new tool and teaching others).
% TODO: Connect to course concepts of funds of knowledge and
%       asset-based views of children and families.

% -----------------------------------------------------------------------
% SECTION 5 — Lessons for My Future Practice
% -----------------------------------------------------------------------
\section{Lessons for My Future Practice}

% TODO: Identify 2-3 concrete commitments for your future classroom.
% TODO: Discuss treating behavior labels as starting points for inquiry,
%       not fixed truths.
% TODO: Explain how you will design learning experiences that surface
%       student strengths early, and how you will self-monitor for
%       deficit thinking.
% TODO: Connect explicitly to course readings:
%       DisCrit, media representation, tragedy vs. non-tragedy models
%       of disability, race and schooling, etc.
% TODO: Name concrete practices: how you will talk about students
%       in meetings, use transcripts/field notes to monitor your language,
%       and engage colleagues around labels.

% -----------------------------------------------------------------------
% SECTION 6 — Conclusion: How My Lens Shifted
% -----------------------------------------------------------------------
\section{Conclusion: How My Lens Shifted}

% TODO: Return to the original stereotypes and summarize how interactions
%       and transcripts complicated those assumptions.
% TODO: Acknowledge the ongoing tension of working within a system that
%       continues to label and track students in ways that can reinforce
%       prejudice.
% TODO: Briefly state how this experience will shape future school entry
%       and connect to longer-term PhD/EQUITAS commitments.

% -----------------------------------------------------------------------
% REFERENCES
% -----------------------------------------------------------------------
\bibliographystyle{apacite}
\bibliography{references}   % create references.bib in Overleaf

\end{document}
